\documentclass[11pt]{article}
\usepackage[margin=1in]{geometry}
\usepackage{amsmath,amssymb}
\usepackage{graphicx}
\usepackage[hidelinks]{hyperref}

\title{Literature-Implied Answer To Problem 4.1 Of arXiv:2107.04662}
\author{Run fa\_banach\_001}
\date{2026-06-28}

\begin{document}
\maketitle

\section*{Status}

\textbf{literature\_implied\_answer (full problem).}
This packet records that Problem 4.1 from Kakol--Leiderman
\cite{KL2021} is answered negatively by combining the problem statement
with a no-surjection theorem recalled in Kakol--Kurka--Sliwa
\cite{KKS2025}, attributed there to Leiderman--Levin--Pestov.
The relation is an agent-identified implication; the supporting source
does not explicitly say that it is resolving Problem 4.1.

\section*{Source Question}

Kakol--Leiderman ask, in Problem 4.1 on PDF page 10:

\begin{center}
\includegraphics[width=0.95\linewidth]{figures/open_problem_crop.png}
\end{center}

In text, the question is whether there exists an infinite compact space
\(X\) admitting a continuous linear surjection
\[
T:C_p(X)\longrightarrow C_k(X)_w.
\]
Here \(C_k(X)_w\) denotes the Banach space \(C(X)\), with its compact-open
norm topology forgotten and replaced by the weak topology.

\section*{Supporting Theorem}

Kakol--Kurka--Sliwa \cite{KKS2025}, in the introduction immediately after
their Problem 1, recall the following restriction from the earlier work
of Leiderman--Levin--Pestov:

\begin{quote}
If \(X\) is a Tychonoff space and \(E\) is an infinite-dimensional normed
space, then there exists no sequentially continuous linear surjection from
\(C_p(X)\) onto \(E_w\).
\end{quote}

\section*{Identification}

Assume that \(X\) is an infinite compact Hausdorff space and that
\[
T:C_p(X)\longrightarrow C_k(X)_w
\]
is a continuous linear surjection.  Put \(E=C(X)\), regarded as a normed
space with the usual supremum norm.  Since \(X\) is infinite compact
Hausdorff, \(C(X)\) is infinite-dimensional: for every finite set of
distinct points of \(X\), normality gives continuous functions separating
those points, so \(C(X)\) has arbitrarily large finite-dimensional
subspaces.

The codomain \(C_k(X)_w\) is precisely \(E_w\).  Moreover, every continuous
map is sequentially continuous.  Thus \(T\) is a sequentially continuous
linear surjection from \(C_p(X)\) onto \(E_w\), contradicting the quoted
Leiderman--Levin--Pestov restriction.  Hence no such infinite compact
space \(X\) exists.

\section*{Scope And Novelty Check}

This is not claimed as a new theorem of the run.  It is a full negative
answer to the extracted Problem 4.1 by direct identification with a known
no-surjection theorem.  The source paper itself proves nonexistence for
infinite metrizable compact \(C\)-spaces; the literature implication above
rules out all infinite compact \(X\).

Cheap run indexes were searched for `2107.04662', `C\_k(X)\_w',
`linear surjection', and related terms, with no exact prior packet found.
Local source search found the later arXiv paper `2509.08634' quoting the
decisive restriction, and web search found the same two primary arXiv
sources.

\begin{thebibliography}{9}
\bibitem{KL2021}
J. Kakol and A. Leiderman,
\emph{On linear continuous operators between distinguished spaces \(C_p(X)\)},
arXiv:2107.04662, 2021.

\bibitem{KKS2025}
J. Kakol, O. Kurka, and W. Sliwa,
\emph{Lotz-Peck-Porta and Rosenthal's theorems for spaces \(C_p(X)\)},
arXiv:2509.08634, 2025.
\end{thebibliography}

\end{document}
